% Options for packages loaded elsewhere
\PassOptionsToPackage{unicode}{hyperref}
\PassOptionsToPackage{hyphens}{url}
\documentclass[
]{article}
\usepackage{xcolor}
\usepackage[margin=1in]{geometry}
\usepackage{amsmath,amssymb}
\setcounter{secnumdepth}{-\maxdimen} % remove section numbering
\usepackage{iftex}
\ifPDFTeX
  \usepackage[T1]{fontenc}
  \usepackage[utf8]{inputenc}
  \usepackage{textcomp} % provide euro and other symbols
\else % if luatex or xetex
  \usepackage{unicode-math} % this also loads fontspec
  \defaultfontfeatures{Scale=MatchLowercase}
  \defaultfontfeatures[\rmfamily]{Ligatures=TeX,Scale=1}
\fi
\usepackage{lmodern}
\ifPDFTeX\else
  % xetex/luatex font selection
\fi
% Use upquote if available, for straight quotes in verbatim environments
\IfFileExists{upquote.sty}{\usepackage{upquote}}{}
\IfFileExists{microtype.sty}{% use microtype if available
  \usepackage[]{microtype}
  \UseMicrotypeSet[protrusion]{basicmath} % disable protrusion for tt fonts
}{}
\makeatletter
\@ifundefined{KOMAClassName}{% if non-KOMA class
  \IfFileExists{parskip.sty}{%
    \usepackage{parskip}
  }{% else
    \setlength{\parindent}{0pt}
    \setlength{\parskip}{6pt plus 2pt minus 1pt}}
}{% if KOMA class
  \KOMAoptions{parskip=half}}
\makeatother
\usepackage{color}
\usepackage{fancyvrb}
\newcommand{\VerbBar}{|}
\newcommand{\VERB}{\Verb[commandchars=\\\{\}]}
\DefineVerbatimEnvironment{Highlighting}{Verbatim}{commandchars=\\\{\}}
% Add ',fontsize=\small' for more characters per line
\usepackage{framed}
\definecolor{shadecolor}{RGB}{248,248,248}
\newenvironment{Shaded}{\begin{snugshade}}{\end{snugshade}}
\newcommand{\AlertTok}[1]{\textcolor[rgb]{0.94,0.16,0.16}{#1}}
\newcommand{\AnnotationTok}[1]{\textcolor[rgb]{0.56,0.35,0.01}{\textbf{\textit{#1}}}}
\newcommand{\AttributeTok}[1]{\textcolor[rgb]{0.13,0.29,0.53}{#1}}
\newcommand{\BaseNTok}[1]{\textcolor[rgb]{0.00,0.00,0.81}{#1}}
\newcommand{\BuiltInTok}[1]{#1}
\newcommand{\CharTok}[1]{\textcolor[rgb]{0.31,0.60,0.02}{#1}}
\newcommand{\CommentTok}[1]{\textcolor[rgb]{0.56,0.35,0.01}{\textit{#1}}}
\newcommand{\CommentVarTok}[1]{\textcolor[rgb]{0.56,0.35,0.01}{\textbf{\textit{#1}}}}
\newcommand{\ConstantTok}[1]{\textcolor[rgb]{0.56,0.35,0.01}{#1}}
\newcommand{\ControlFlowTok}[1]{\textcolor[rgb]{0.13,0.29,0.53}{\textbf{#1}}}
\newcommand{\DataTypeTok}[1]{\textcolor[rgb]{0.13,0.29,0.53}{#1}}
\newcommand{\DecValTok}[1]{\textcolor[rgb]{0.00,0.00,0.81}{#1}}
\newcommand{\DocumentationTok}[1]{\textcolor[rgb]{0.56,0.35,0.01}{\textbf{\textit{#1}}}}
\newcommand{\ErrorTok}[1]{\textcolor[rgb]{0.64,0.00,0.00}{\textbf{#1}}}
\newcommand{\ExtensionTok}[1]{#1}
\newcommand{\FloatTok}[1]{\textcolor[rgb]{0.00,0.00,0.81}{#1}}
\newcommand{\FunctionTok}[1]{\textcolor[rgb]{0.13,0.29,0.53}{\textbf{#1}}}
\newcommand{\ImportTok}[1]{#1}
\newcommand{\InformationTok}[1]{\textcolor[rgb]{0.56,0.35,0.01}{\textbf{\textit{#1}}}}
\newcommand{\KeywordTok}[1]{\textcolor[rgb]{0.13,0.29,0.53}{\textbf{#1}}}
\newcommand{\NormalTok}[1]{#1}
\newcommand{\OperatorTok}[1]{\textcolor[rgb]{0.81,0.36,0.00}{\textbf{#1}}}
\newcommand{\OtherTok}[1]{\textcolor[rgb]{0.56,0.35,0.01}{#1}}
\newcommand{\PreprocessorTok}[1]{\textcolor[rgb]{0.56,0.35,0.01}{\textit{#1}}}
\newcommand{\RegionMarkerTok}[1]{#1}
\newcommand{\SpecialCharTok}[1]{\textcolor[rgb]{0.81,0.36,0.00}{\textbf{#1}}}
\newcommand{\SpecialStringTok}[1]{\textcolor[rgb]{0.31,0.60,0.02}{#1}}
\newcommand{\StringTok}[1]{\textcolor[rgb]{0.31,0.60,0.02}{#1}}
\newcommand{\VariableTok}[1]{\textcolor[rgb]{0.00,0.00,0.00}{#1}}
\newcommand{\VerbatimStringTok}[1]{\textcolor[rgb]{0.31,0.60,0.02}{#1}}
\newcommand{\WarningTok}[1]{\textcolor[rgb]{0.56,0.35,0.01}{\textbf{\textit{#1}}}}
\usepackage{graphicx}
\makeatletter
\newsavebox\pandoc@box
\newcommand*\pandocbounded[1]{% scales image to fit in text height/width
  \sbox\pandoc@box{#1}%
  \Gscale@div\@tempa{\textheight}{\dimexpr\ht\pandoc@box+\dp\pandoc@box\relax}%
  \Gscale@div\@tempb{\linewidth}{\wd\pandoc@box}%
  \ifdim\@tempb\p@<\@tempa\p@\let\@tempa\@tempb\fi% select the smaller of both
  \ifdim\@tempa\p@<\p@\scalebox{\@tempa}{\usebox\pandoc@box}%
  \else\usebox{\pandoc@box}%
  \fi%
}
% Set default figure placement to htbp
\def\fps@figure{htbp}
\makeatother
\setlength{\emergencystretch}{3em} % prevent overfull lines
\providecommand{\tightlist}{%
  \setlength{\itemsep}{0pt}\setlength{\parskip}{0pt}}
\usepackage{bookmark}
\IfFileExists{xurl.sty}{\usepackage{xurl}}{} % add URL line breaks if available
\urlstyle{same}
\hypersetup{
  pdftitle={Análise Exploratória de Dados - Dados Turísticos},
  pdfauthor={Gabriel Rodrigues},
  hidelinks,
  pdfcreator={LaTeX via pandoc}}

\title{Análise Exploratória de Dados - Dados Turísticos}
\author{Gabriel Rodrigues}
\date{2026-05-01}

\begin{document}
\maketitle

\begin{Shaded}
\begin{Highlighting}[]
\FunctionTok{setwd}\NormalTok{(}\StringTok{\textquotesingle{}/home/gabrielr/Documentos/UFMG/2026.1/ICD/Prática/Projeto/Data{-}Science{-}Project/\textquotesingle{}}\NormalTok{)}
\end{Highlighting}
\end{Shaded}

\section{Dados Macroeconômicos Turísticos
Mundiais}\label{dados-macroeconuxf4micos-turuxedsticos-mundiais}

Primeiramente, vamos importar os dados, tratá-los e visualizar alguns
pontos chave.

\begin{Shaded}
\begin{Highlighting}[]
\CommentTok{\# Carregando as bibliotecas principais}
\FunctionTok{library}\NormalTok{(tidyverse)}
\end{Highlighting}
\end{Shaded}

\begin{verbatim}
## -- Attaching core tidyverse packages ------------------------ tidyverse 2.0.0 --
## v dplyr     1.2.1     v readr     2.2.0
## v forcats   1.0.1     v stringr   1.6.0
## v ggplot2   4.0.3     v tibble    3.3.1
## v lubridate 1.9.5     v tidyr     1.3.2
## v purrr     1.2.2     
## -- Conflicts ------------------------------------------ tidyverse_conflicts() --
## x dplyr::filter() masks stats::filter()
## x dplyr::lag()    masks stats::lag()
## i Use the conflicted package (<http://conflicted.r-lib.org/>) to force all conflicts to become errors
\end{verbatim}

\begin{Shaded}
\begin{Highlighting}[]
\CommentTok{\# Lendo os dados brutos}

\NormalTok{df\_turismo }\OtherTok{\textless{}{-}} \FunctionTok{read.csv}\NormalTok{(}\StringTok{"Tourism Quality/Tourism Quality Data.csv"}\NormalTok{)}

\CommentTok{\# Ajuste inicial de tipos de dados}
\NormalTok{df\_turismo }\OtherTok{\textless{}{-}}\NormalTok{ df\_turismo }\SpecialCharTok{\%\textgreater{}\%}
  \FunctionTok{mutate}\NormalTok{(}
    \AttributeTok{Ano =} \FunctionTok{as.factor}\NormalTok{(Ano), }\CommentTok{\# Transformando o Ano em fator (categoria)}
    \AttributeTok{IncomeGroup =} \FunctionTok{factor}\NormalTok{(IncomeGroup, }
                         \AttributeTok{levels =} \FunctionTok{c}\NormalTok{(}\StringTok{"Low income"}\NormalTok{, }\StringTok{"Lower middle income"}\NormalTok{, }
                                    \StringTok{"Upper middle income"}\NormalTok{, }\StringTok{"High income"}\NormalTok{)) }\CommentTok{\# Ordenando a renda}
\NormalTok{  )}

\CommentTok{\# Visualização inicial da base final}
\FunctionTok{glimpse}\NormalTok{(df\_turismo)}
\end{Highlighting}
\end{Shaded}

\begin{verbatim}
## Rows: 333
## Columns: 24
## $ Sigla                                                      <chr> "AGO", "AGO~
## $ País                                                       <chr> "Angola", "~
## $ Ano                                                        <fct> 2019, 2021,~
## $ Region                                                     <chr> "Sub-Sahara~
## $ IncomeGroup                                                <fct> Lower middl~
## $ Continente                                                 <chr> "AF", "AF",~
## $ Cultural.Demand                                            <dbl> 1.138445, 1~
## $ Cultural.Resources                                         <dbl> 1.209611, 1~
## $ Geographic.Dispersion                                      <dbl> 2.607143, 2~
## $ Natural.Resources                                          <dbl> 3.003402, 3~
## $ Natural.Demand                                             <dbl> 1.057587, 1~
## $ Non.Leisure.Resources                                      <dbl> 1.270534, 1~
## $ Non.Leisure.Demand                                         <dbl> 1.052973, 1~
## $ Openness.to.Tourism                                        <dbl> 2.199659, 2~
## $ Prioritization.of.Tourism                                  <dbl> 3.618593, 3~
## $ Air.Transport.Infrastructure                               <dbl> 2.226030, 2~
## $ Ground.and.Port.Infrastructure                             <dbl> 1.816664, 1~
## $ T.T.Employment.Multiplier                                  <dbl> 4.756657, 5~
## $ T.T.GDP.Multiplier                                         <dbl> 4.680924, 5~
## $ T.T.Government.Expenditure.....of.Government.Budget        <dbl> 1.672798, 1~
## $ T.T.High.Wage.Jobs......of.Employment.in.High.Wage.Sectors <dbl> 6.075938, 5~
## $ Domestic.T.T.spending.....of.Internal.T.T.Spending         <dbl> 64.45112, 8~
## $ T.T.industry.Share.of.GDP.....of.total.GDP                 <dbl> 1.3524970, ~
## $ T.T.industry.GDP                                           <dbl> 1328.0020, ~
\end{verbatim}

\begin{Shaded}
\begin{Highlighting}[]
\FunctionTok{unique}\NormalTok{(df\_turismo}\SpecialCharTok{$}\NormalTok{Ano)}
\end{Highlighting}
\end{Shaded}

\begin{verbatim}
## [1] 2019 2021 2024
## Levels: 2019 2021 2024
\end{verbatim}

\subsection{Notas metodológicas acerca da da natureza dos
dados}\label{notas-metodoluxf3gicas-acerca-da-da-natureza-dos-dados}

Antes de adentrarmos nas análises exploratórias, é fundamental tecer
comentários sobre a natureza estrutural deste dataset (Tourism Quality
Data) e, consequentemente, sobre os limites das inferências estatísticas
que podem ser extraídas dele.

\begin{itemize}
\tightlist
\item
  \textbf{Unidade de análise Macro (País-Ano)}
\end{itemize}

Os dados trabalhados consistem em um painel macroeconômico. Nossa
observação mínima não é o indivíduo, o viajante, ou empresas (dados
analíticos), mas a entidade ``País'' ou ``Território'' em recortes
temporais específicos (2019, 2021 e 2024) sob determinados critérios.

\begin{itemize}
\tightlist
\item
  \textbf{Microdados vs Macrodados}
\end{itemize}

Não é possível utilizar esta base para análises probabilísticas de
comportamento individual (ex.: ``qual a probabilidade de um turista
gastar mais em um destino com mais recursos naturais?''). O dado
individual foi sumarizado em indicadores consolidados pelo macro do
país, como multiplicadores de emprego, demanda natural, dispersão
geográfica etc.

\begin{itemize}
\tightlist
\item
  \textbf{Índices e Scores Padronizados}
\end{itemize}

A maioria das variáveis qualitativas presentes no dataset (Tourism
Quality Data) já são estatísticas consolidadas sob a forma de scores
numéricos.

\begin{itemize}
\tightlist
\item
  \textbf{Abordagem estatística adotada}
\end{itemize}

Cientes de tais premissas, as análises descritivas e inferenciais
focarão em padrões sistêmicos. Utilizaremos matrizes de correlaão,
análise de variância entre estados globais e modelos de regressão para
entender como a infraestrutura e a priorização governamental se
relacionam no macro com a geração de riqueza do setor de turismo.

\subsubsection{O que os dados nos permitem responder de nossas perguntas
iniciais?}\label{o-que-os-dados-nos-permitem-responder-de-nossas-perguntas-iniciais}

A estrutura atual do dataset é rica em variáveis de corte transversal
(cross-section) e agregações categóricas. Portanto, a base é
perfeitamente adequada para responder a perguntas focadas em perfis,
correlações, agrupamentos e comparações estruturais.

Os dados suportam com alto rigor estatístico as seguintes investigações
propostas:

\begin{itemize}
\item
  \textbf{Comparações e Perfis Estruturais:} Questões sobre Diferenças
  regionais (África Subsaariana vs.~Zona do Euro vs.~Mercosul
  vs.~BRICS\ldots; França vs Inglaterra vs Rússia vs.~Brasil), Impacto
  do IncomeGroup, e Turismo de lazer vs.~negócios. Podemos utilizar
  Análise de Variância (ANOVA) e boxplots para provar se a renda ou a
  região definem categorias com comportamentos estatisticamente
  distintos.
\item
  \textbf{Eficiência e Relações de Recursos:} Questões sobre Recursos
  vs.~Demanda, Eficiência Econômica e Recursos Naturais vs.~Culturais.
  Ao criarmos razões (ex: PIB / Recursos Naturais), podemos facilmente
  mapear quais países ``fazem mais com menos'' e identificar padrões de
  especialização (matrizes de correlação).
\item
  \textbf{Modelagem e Agrupamentos:}
\end{itemize}

Além das análises descritivas e comparativas, também podemos utilizar
técnicas simples de modelagem e agrupamento para explorar melhor os
padrões presentes nos dados.

\begin{itemize}
\tightlist
\item
  \textbf{Relações entre variáveis:} Podemos investigar como diferentes
  indicadores se relacionam entre si. Por exemplo, é possível verificar
  se países com melhor infraestrutura turística tendem a apresentar
  maior receita ou demanda. Para isso, utilizamos ferramentas como a
  matriz de correlação e gráficos de dispersão.
\item
  \textbf{Modelos de regressão (introdução):} Uma abordagem inicial é
  utilizar regressões lineares simples para avaliar se uma variável pode
  ajudar a explicar outra. Por exemplo, podemos testar se indicadores
  como investimento em turismo ou qualidade de serviços estão associados
  ao desempenho econômico do setor. É importante destacar que esses
  modelos indicam associação, não causalidade.
\item
  \textbf{Agrupamento de países (clusters):} Também podemos aplicar
  técnicas de agrupamento, como o K-means, para identificar grupos de
  países com características semelhantes. Esses grupos podem revelar
  padrões interessantes, como países com alto investimento e baixa
  eficiência, ou países que obtêm bons resultados com menos recursos.
\item
  Identificação de padrões e outliers: Essas abordagens permitem
  identificar tanto padrões gerais quanto casos atípicos (outliers), que
  podem representar países com desempenho muito acima ou abaixo do
  esperado em relação aos demais.
\end{itemize}

De forma geral, essas técnicas ajudam a organizar melhor a informação
contida nos dados e a gerar hipóteses sobre o funcionamento do setor
turístico em nível global.

\subsubsection{E o que eles não nos permitem
responder?}\label{e-o-que-eles-nuxe3o-nos-permitem-responder}

Algumas das perguntas propostas exigem estruturas de dados que não
possuímos (como séries temporais longas ou testes de causalidade
controlados). Tentar respondê-las diretamente com esta base resultará em
conclusões estatisticamente frágeis. As principais restrições recaem
sobre:

\begin{itemize}
\item
  \textbf{Tendências Históricas e Mudanças Estruturais de Longo Prazo:}
  Perguntas sobre Evolução de demandas ao longo do tempo e Países
  mudando de perfil.

  \begin{itemize}
  \tightlist
  \item
    O porquê: possuímos apenas três recortes temporais (2019, 2021 e
    2024). Em estatística, três pontos não configuram uma ``tendência
    histórica''. O que os dados de fato mostram é o choque agudo da
    pandemia de COVID-19 (a queda em 2021 e a recuperação em 2024). Não
    é possível afirmar que um país sofreu uma ``mudança estrutural de
    modelo de turismo''; ele apenas pode ter tido uma recuperação mais
    lenta ou rápida da crise sanitária.
  \end{itemize}
\item
  \textbf{Causalidade e Trade-offs Exclusivos:} Perguntas como
  Infraestrutura como fator limitante para crescimento e Trade-offs
  (investir em X reduz Y).

  \begin{itemize}
  \tightlist
  \item
    O porquê: correlação não implica causalidade. A base nos permite
    afirmar que ``países com maior infraestrutura tendem a ter maior
    receita'', mas não prova que o investimento financeiro direto na
    infraestrutura causou a receita. Da mesma forma, encontrar uma
    correlação negativa entre duas variáveis não prova a existência de
    um trade-off consciente de políticas públicas (onde o governo
    escolheu tirar dinheiro de X para colocar em Y).
  \end{itemize}
\item
  \textbf{Convergência de Longo Prazo:} Pergunta sobre Países mais
  abertos convergem para níveis mais altos.

  \begin{itemize}
  \tightlist
  \item
    O porquê: novamente, a janela de 5 anos (atravessada por uma
    pandemia que fechou fronteiras artificialmente) impede uma análise
    de convergência real. Podemos, no máximo, avaliar a elasticidade da
    recuperação de países abertos versus fechados no pós-crise.
  \end{itemize}
\end{itemize}

\subsection{Análises e plotagens}\label{anuxe1lises-e-plotagens}

\subsubsection{Matriz de correlação}\label{matriz-de-correlauxe7uxe3o}

\end{document}
